% Generated draft: requires mathematical and source review.
\documentclass[11pt]{article}
\usepackage[margin=25mm]{geometry}
\usepackage{fontspec}
\setmainfont{Latin Modern Roman}
\usepackage{amsmath,amssymb,mathtools,unicode-math}
\setmathfont{Latin Modern Math}
\usepackage{enumitem,needspace,xurl,hyperref}
\setlength{\parindent}{0pt}
\setlength{\parskip}{5pt}
\setlength{\emergencystretch}{3em}
\newcommand{\N}{\mathbb{N}}
\newcommand{\Z}{\mathbb{Z}}
\newcommand{\Q}{\mathbb{Q}}
\newcommand{\R}{\mathbb{R}}
\newcommand{\C}{\mathbb{C}}
\newcommand{\F}{\mathbb{F}}
\newcommand{\A}{\mathbb{A}}
\newcommand{\PP}{\mathbb{P}}
\newcommand{\E}{\mathbb{E}}
\newcommand{\eps}{\varepsilon}
\newcommand{\dd}{\,\mathrm{d}}
\begin{document}
\section*{UnsolvedMath ID: 20003088}
\textbf{Source problem code:} catalogue entry
\subsection{Definitions and mathematical statement}
The extraction-damaged source asks whether a Thurston unit sphere at a point of Teichm\"uller space determines global stretch lines or geodesics. The supplied interpretation retains the full linearly embedded cotangent dual co-sphere, rather than only its abstract topology or face incidence. For an oriented finite-type surface $S$ with negative Euler characteristic, write $\mathcal T(S)$ for marked hyperbolic Teichm\"uller space and $\ell_\lambda(X)$ for measured-lamination length. The asymmetric tangent norm and distance are
\[
\|v\|_{\mathrm{Th},X}=\sup_{\lambda\ne0}\dd\log\ell_\lambda|_X(v),\qquad
 d_{\mathrm{Th}}(X,Y)=\log\sup_{\lambda\ne0}\frac{\ell_\lambda(Y)}{\ell_\lambda(X)}.
\]
Let $K_X\subset T_X^*\mathcal T(S)$ be the closed dual unit ball, characterized by its support function $\sup_{p\in K_X}p(v)=\|v\|_{\mathrm{Th},X}$, and let $\partial K_X$ be its co-sphere. Asymmetry is retained; it is not replaced by a symmetric norm.

The aid's rigidity reduction assumes a real linear isomorphism
\[
A:T_Y^*\mathcal T(S)\longrightarrow T_X^*\mathcal T(S),\qquad
A(\partial K_Y)=\partial K_X.
\]
Convexity determines the balls from their boundaries, and polarity gives $\|A^*v\|_{\mathrm{Th},Y}=\|v\|_{\mathrm{Th},X}$ for $A^*:T_X\mathcal T(S)\to T_Y\mathcal T(S)$. In the surface range of the infinitesimal rigidity theorem invoked by the aid, this tangent isometry is induced by an extended mapping class $f$ taking $X$ to $Y$. It follows that the pointed global Thurston metric is determined up to that action. The source does not supply the theorem's full surface range or low-complexity exceptions, so those are required hypotheses rather than silently universal assertions.

The resulting global isometry transports whole families of paths, not a uniquely selected geodesic. For the standard stretch path based at $X$ along a specified complete geodesic lamination $\Lambda$, the reported equivariance is
\[
f(\operatorname{Stretch}_t(X,\Lambda))
=\operatorname{Stretch}_t(fX,f\Lambda).
\]
It also preserves oriented geodesic envelopes, expressed by
\[
\operatorname{Env}(X,Z)=\{W:d_{\mathrm{Th}}(X,W)+d_{\mathrm{Th}}(W,Z)=d_{\mathrm{Th}}(X,Z)\}.
\]
Thus full embedded co-sphere data provide a sufficient rigidity-level formulation of the source question, without an explicit reconstruction algorithm.

For closed $S$ of genus at least two, the aid further reports that extreme points of the polar tangent unit body recognize maximal chain-recurrent stretch directions. It also reports basepoint independence of abstract tangent-face combinatorics, showing that face incidence alone cannot retain the affine-linear data needed for this conclusion. Neither precise extreme-point correspondence nor the face-combinatorics theorem is stated or cited fully in the record. The source's remark about infinitesimal lamination-length behavior giving global coordinates is retained as background motivation, not used to identify abstract face data with the full embedded convex body.
\subsection{Short English statement}
The supplied precise answer keeps the entire cotangent co-sphere inside its linear space. Polarity turns a linear equivalence of such co-spheres into a tangent Thurston isometry, and the applicable rigidity theorem then gives a global mapping-class isometry transporting stretch paths and geodesic envelopes. Abstract face combinatorics is weaker and is reported insufficient. Surface-range and normalization details still need the missing references.
\subsection{Sources}
\begin{itemize}
\item Local registry: permanent catalogue entry, exact source code catalogue entry; \texttt{aim-topology-notes.json}, \texttt{notes[175]}.
\item AIM workshop \emph{Lipschitz metric on Teichmueller space}, Problem 1.3, attributed in the extraction to ``F. Gu\'ertitaud'', \url{https://aimath.org/pastworkshops/lipschitzteichproblems.pdf}. The source names Thurston's infinitesimal measured-lamination-length coordinates as motivation. No complete publication citations are supplied.
\item Supplied research attempt 1, registry updated 14 August 2026: partial result, low novelty confidence, machine-generated and not peer reviewed. It invokes infinitesimal rigidity, the polar extreme-point/stretch-direction result for closed genus at least two, and recent basepoint independence of tangent face combinatorics without titles, authors, theorem numbers, or complete authoritative references.
\item Statement status is \texttt{reconstructed\_unverified}; clean statement is null. Verification explicitly retains the canonical extraction errors rather than supplying a checked original formula.
\end{itemize}
\end{document}
